\chapter{Estado del arte}\label{cap:estado-arte}

\section{Plataformas y servicios de \emph{fact-checking} asistido}

Existen diversas iniciativas, tanto académicas como comerciales, que
abordan total o parcialmente el problema:

\begin{itemize}[leftmargin=*,itemsep=0.2em]
  \item \textbf{ClaimBuster} \cite{hassan2017claimbuster}: pipeline
        académico que identifica afirmaciones susceptibles de
        verificación.
  \item \textbf{Full Fact} y \textbf{Chequeado}: organizaciones de
        verificación humana que han incorporado herramientas asistidas
        para priorizar contenido a revisar.
  \item \textbf{Google Fact Check Tools} y \textbf{ClaimReview}:
        infraestructura abierta para indexar verificaciones publicadas
        por terceros.
\end{itemize}

Frente a estas, el presente trabajo persigue un sistema integrado
\emph{end-to-end}: el usuario introduce un texto, el sistema lo
procesa y devuelve un veredicto con evidencias. La parte de
\emph{retrieval} se delega a una API de búsqueda externa para evitar
mantener un índice propio de la web.

\section{Tecnologías candidatas y decisiones}

\subsection{Backend}

\paragraph{Python con FastAPI.} Se elige FastAPI por su
\emph{performance} (basado en Starlette y Pydantic), tipado estático
y generación automática de OpenAPI. La alternativa Flask se descarta
por carecer de validación tipada nativa; Django REST Framework por
imponer una capa ORM no necesaria, dado que se usa Supabase.

\subsection{Frontend}

\paragraph{React + Vite + Tailwind + shadcn/ui.} React se selecciona
por madurez y ecosistema. Vite ofrece tiempos de arranque y \emph{hot
reload} significativamente mejores que Create React App. Tailwind
permite construir un sistema de diseño consistente sin abandonar el
flujo \emph{utility-first}. shadcn/ui aporta componentes accesibles
sobre Radix UI sin acoplar a un \emph{theme provider} cerrado.

\paragraph{Estado: Zustand.} Frente a Redux (boilerplate) o Context API
(rendimiento), Zustand ofrece una API minimalista, integración
inmediata con React 18 y \emph{stores} compuestos.

\subsection{Backend como servicio: Supabase}

Supabase ofrece PostgreSQL gestionado, Auth (incluyendo OAuth con
Google) y \emph{storage}, con Row Level Security activable por tabla.
La alternativa Firebase se descarta por su modelo NoSQL menos adecuado
para datos relacionales (cuentas, suscripciones, historial).

\subsection{Pagos: Stripe}

Stripe es el estándar \emph{de facto} para SaaS por API estable,
documentación, integración con \emph{webhooks} y soporte de
suscripciones recurrentes con periodos de prueba.

\subsection{Extensión de navegador}

Manifest V3 es el formato vigente en Chromium y Firefox; impone uso de
\emph{service workers} y restricciones más estrictas de CSP. La
implementación reutiliza el mismo cliente de Supabase y la misma API
REST que el frontend principal.

\section{Conclusión}

La combinación FastAPI + React + Supabase + Stripe + Docker Compose
representa un \emph{stack} probado en SaaS modernos y permite cumplir
los objetivos del proyecto sin acoplar el sistema a un proveedor
\emph{cloud} concreto.
